% ============================================================================
% CHAPTER 3 - METHODS / PARAMETERS
% File: docs/thesis/methods/parameters/group2_core_costs.tex
% Group 2: core costs (Stage F6). OQ-016, OQ-017, OQ-018, OQ-022, OQ-054.
% Theory: Section sec:lcoe_theory in ch2_renewable_potential_foundations.tex
% Values must match config/parameters.json. References marked % verify are unchecked.
% ============================================================================

\section{Core Cost Parameters (Stage F6)}
\label{sec:params_group2}

\subsection{Scope and Conventions}
\label{subsec:g2_scope}

This section fixes the capital cost, fixed operating cost, discount rate,
lifetime and reference commissioning year used in the LCOE of Stage F6
(Equation~\ref{eq:theory_lcoe}). Grid, substation and road costs and
degradation belong to Group 3.

Reference year and currency. The technology and cost vintage is
$t_{\mathrm{ref}}=2024$ (OQ-018), the commissioning cohort of IRENA's
cost report, which supplies the capital costs, the operating costs and
the discount rates of this section \cite{irena2025costs2024}. All values
are in constant 2024 USD and the discount rates are real, so no deflator
and no nominal-to-real conversion is needed (OQ-054).

Electrical basis. IRENA reports PV total installed costs per
kW$_{\mathrm{dc}}$ \cite{irena2025costs2024}, the basis used by GeoFREA,
so no conversion is applied. Stage F5 does not model AC output.

Sampling. Parameters sampled in Stage F6 (U-03) use uniform
distributions. The central value is the one used in deterministic runs.

\subsection{Capital Cost (OQ-017)}
\label{subsec:g2_capex}

Definition: total installed cost per unit of capacity at commissioning.
Unit: USD$_{2024}$/kW$_{\mathrm{dc}}$ (PV), USD$_{2024}$/kW (wind).
Enters: Equation~\ref{eq:theory_lcoe}; sampled in Stage F6.

IRENA reports, for each country and technology, the capacity-weighted
average and the 5th and 95th percentiles of project costs
\cite{irena2025costs2024}. Brazil and India have country values;
Portugal uses the European regional values. A regional value of the same
quantity and technology is a Tier 2 transfer, not a proxy.

The project distribution is right-skewed, so a uniform distribution on
the 5th--95th percentile interval would have a mean well above the
reported average and would raise every sampled LCOE. The sampling
interval is instead centred on the weighted average, with half-width
equal to the distance between the average and the 5th percentile:
\begin{equation}
    \mathrm{CAPEX} \sim \mathcal{U}\!\left[P_5,\; 2\bar{c}-P_5\right],
    \label{eq:g2_capex_interval}
\end{equation}
where $\bar{c}$ is the weighted average. The mean of the draws equals
the reported average. The upper tail beyond $2\bar{c}-P_5$ is excluded;
it reflects small, remote or complex projects, part of which GeoFREA
represents explicitly through its grid and road costs (Group 3).

\begin{table}[htbp]
    \centering
    \caption{Capital cost (2024 USD, 2024 commissioning cohort).}
    \label{tab:g2_capex}
    \small
    \begin{tabular}{llrrr}
        \toprule
        Technology & Country & IRENA $\bar{c}$ (central) & IRENA $P_5$ / $P_{95}$ & Sampling range \\
        \midrule
        PV (USD/kW$_{\mathrm{dc}}$) & BRA & 724 & 503 / 1152 & 503--945 \\
                                    & IND & 525 & 415 / 831  & 415--635 \\
                                    & PRT (Europe) & 779 & 524 / 1312 & 524--1034 \\
        \midrule
        Wind (USD/kW)               & BRA & 1307 & 1015 / 1628 & 1015--1599 \\
                                    & IND & 1110 & 984 / 1270  & 984--1236 \\
                                    & PRT (Europe) & 1659 & 1211 / 2232 & 1211--2107 \\
        \bottomrule
    \end{tabular}
\end{table}

Grid connection. For PV, IRENA's total installed cost includes the
plant-side connection (substation, transformers and medium-voltage
cabling) but does not state whether it includes a line to a distant
transmission node \cite{irena2025costs2024}. For onshore wind the report
does not state the scope. GeoFREA's connection costs (Group 3, OQ-001)
must therefore exclude the plant-side substation for PV, and the wind
case is resolved in Group 3.

Source: \cite{irena2025costs2024}. Tier 1 (BRA, IND), Tier 2 (PRT).
Distribution: uniform. Proxy: no. Price year: 2024.

\subsection{Fixed Operating Cost (OQ-016)}
\label{subsec:g2_opex}

Definition: annual fixed operation and maintenance cost as a fraction of
the capital cost (opex\_fixed\_frac). Unit: 1/yr. Enters:
Equation~\ref{eq:theory_lcoe} as $F=f\cdot I$; sampled in Stage F6.
Variable operating cost is set to zero, following the all-in fixed
convention of the cost source.

Source values in USD/kW/yr are divided by the weighted-average capital
cost of the same country and technology, on the same capacity basis. For
PV in Brazil and India the working range is 6.99--11.0~USD/kW/yr,
extended to the IRENA global all-in value of 13.1~USD/kW/yr to cover the
difference between contract prices and all-in costs
\cite{irena2025costs2024}. For wind, single country values (Brazil 29,
India 23~USD/kW/yr) receive a $\pm20\%$ range by author judgment. The
IRENA O\&M table has no Portuguese wind row; Portugal uses the European
rows of that table, which range from 29 (Spain) to 53~USD/kW/yr, with the
midpoint as central value.

\begin{table}[htbp]
    \centering
    \caption{Fixed O\&M. Fractions are relative to the IRENA weighted
    average capital cost.}
    \label{tab:g2_opex}
    \small
    \begin{tabular}{llrrrc}
        \toprule
        Tech. & Country & USD/kW/yr (central; range) & Central $f$ & Range $f$ & Tier (value / range) \\
        \midrule
        PV   & BRA & 9.0; 6.99--13.1  & 0.0124 & 0.0097--0.0181 & 2 / 2 \\
             & IND & 9.0; 6.99--13.1  & 0.0171 & 0.0133--0.0250 & 2 / 2 \\
             & PRT & 19.45; 18.2--20.7 & 0.0250 & 0.0234--0.0266 & 2 / 2 \\
        Wind & BRA & 29; 23.2--34.8   & 0.0222 & 0.0178--0.0266 & 1 / 3 \\
             & IND & 23; 18.4--27.6   & 0.0207 & 0.0166--0.0249 & 1 / 3 \\
             & PRT & 41; 29--53       & 0.0247 & 0.0175--0.0319 & 2 / 2 \\
        \bottomrule
    \end{tabular}
\end{table}

Source: \cite{irena2025costs2024}. Distribution: uniform. Proxy: no.
Price year: 2024.

\subsection{Discount Rate (OQ-022)}
\label{subsec:g2_wacc}

Definition: real after-tax weighted average cost of capital. Unit: 1/yr.
Enters: Equation~\ref{eq:theory_lcoe}; sampled in Stage F6.

IRENA's LCOE for the 2024 cohort uses country- and technology-specific
real after-tax WACC values from a benchmark cost-of-capital tool, with a
minimum floor for high-risk countries and default values of 5\% (OECD and
China) and 7.5\% (rest of the world) for countries the benchmark does not
cover; the values reflect conditions prevailing in 2023
\cite{irena2025costs2024}. The country values are published only as an
unlabeled map, whose title states real after-tax rates while its legend
reads nominal; the text and the standardized-assumptions table of the
same report state real rates, which is the reading adopted. The report
also publishes regional values, which GeoFREA uses:

\begin{table}[htbp]
    \centering
    \caption{Real after-tax discount rate (IRENA regional WACC, 2024).}
    \label{tab:g2_wacc}
    \small
    \begin{tabular}{lllccc}
        \toprule
        Country & IRENA region & Tech. & Central & Sampling range & Tier (value / range) \\
        \midrule
        BRA & South America & PV   & 7.7\% & 6.7--8.7\% & 2 / 3 \\
            &               & Wind & 7.7\% & 6.7--8.7\% & 2 / 3 \\
        IND & Asia          & PV   & 4.0\% & 3.0--5.0\% & 2 / 3 \\
            &               & Wind & 3.8\% & 2.8--4.8\% & 2 / 3 \\
        PRT & Europe        & PV   & 4.2\% & 3.2--5.2\% & 2 / 3 \\
            &               & Wind & 3.7\% & 2.7--4.7\% & 2 / 3 \\
        \bottomrule
    \end{tabular}
\end{table}

The sampling range is $\pm1$ percentage point around the regional value,
by author judgment. The Asian value is dominated by China and likely
understates the cost of capital in India. It is consistent, however, with
IRENA's earlier country estimates for India (7.1\% for PV and 8.4\% for
wind, nominal, 2019--2021) once converted with the Reserve Bank of India
inflation target of 4\% (3.0\% and 4.2\% real) \cite{irena2023financing}.
For Brazil, the IRENA country comparison of financing shares places
Brazilian PV at about 8\% \cite{irena2025costs2024}, inside the sampling
range.

Source: \cite{irena2025costs2024}. Distribution: uniform. Proxy: no.

\subsection{Lifetime (OQ-018)}
\label{subsec:g2_lifetime}

Definition: economic lifetime $n$. Unit: yr. Enters:
Equation~\ref{eq:theory_lcoe}; sampled in Stage F6.

Central value 30 years for both technologies, the technical lifetime in
NREL's utility-scale PV and land-based wind assumptions
\cite{nrelATB2024definitions}. Sampled as a discrete uniform integer on
$\{25,\dots,30\}$, each year with probability 1/6; the lower bound equals
the economic life of IRENA's standardized LCOE \cite{irena2025costs2024}.
The central value sits at the upper end of the range, so the uncertainty
is one-sided by design. Because IRENA uses 25 years, GeoFREA's central
LCOE is lower than IRENA's for identical inputs; plausibility comparisons
with IRENA LCOE account for this difference.

Tier 2 (value and range). Proxy: no.

\subsection{Summary}
\label{subsec:g2_summary}

\begin{table}[htbp]
    \centering
    \caption{Group 2 parameters (2024 USD, real).}
    \label{tab:g2_summary}
    \small
    \begin{tabularx}{\textwidth}{@{}>{\raggedright\arraybackslash}p{3.6cm}
        >{\raggedright\arraybackslash}X
        >{\centering\arraybackslash}p{1.2cm}
        >{\centering\arraybackslash}p{1.6cm}@{}}
        \toprule
        Parameter & Value & Sampled & Tier \\
        \midrule
        $t_{\mathrm{ref}}$ & 2024 & no & 1 \\
        CAPEX & Table~\ref{tab:g2_capex} & yes & 1 / 2 \\
        Fixed O\&M fraction & Table~\ref{tab:g2_opex}; variable O\&M = 0 & yes & 1--3 \\
        Discount rate (real) & Table~\ref{tab:g2_wacc} & yes & 2 / 3 \\
        Lifetime & 30 yr, discrete U\{25, \dots, 30\} & yes & 2 \\
        \bottomrule
    \end{tabularx}
\end{table}

\subsection{Limitations}
\label{subsec:g2_limitations}

Portugal uses European regional costs, and all discount rates are
regional averages; the Asian rate is likely low for India. The truncated
upper tail of the capital-cost distribution removes the most expensive
projects, which is consistent only if Group 3 represents the connection
costs that drive them. Some ranges are author judgment ($\pm20\%$ for
single O\&M values, $\pm1$ percentage point for discount rates). Treating
O\&M as a fraction of capital cost makes absolute O\&M scale with sampled
capital cost. Discount rates reflect 2023 financial conditions and are
held constant over the plant life.
